\chapter{Conclusiones}

Este trabajo propuso y evaluó un modelo híbrido de aprendizaje profundo para el pronóstico horario del costo marginal en el Sistema Eléctrico Nacional (SEN), combinando un Temporal Fusion Transformer (TFT) con dos estrategias alternativas de \textit{features} exógenas, Prophet más clima (P+C) y posición solar más clima (S+C), y complementando la predicción puntual con intervalos de Conformal Prediction (CP) calibrados para cuantificar la incertidumbre. Este capítulo resume los hallazgos principales del trabajo, explicita sus contribuciones, reconoce sus limitaciones y propone líneas de continuación.

\section{Síntesis de resultados}

El primer resultado central del trabajo es que \textbf{la estrategia más simple es también la más precisa}: S+C, que reemplaza la descomposición estadística de Prophet y el ensamblaje de Stacking Optimization por la posición solar calculada de forma determinista mediante \texttt{pvlib}, supera a P+C en 14 de las 16 combinaciones de barra y arquitectura evaluadas, diferencia estadísticamente significativa (Diebold-Mariano, Newey-West) en 12 de esas 14, y en 13 de las 16 combinaciones en total (Sección \ref{sec:pc_vs_sc}). El análisis de interpretabilidad de la Variable Selection Network confirma que esta ventaja no es casual: la posición solar domina el decoder del TFT en seis de las ocho barras del SEN, evidencia directa de que el modelo efectivamente aprende a apoyarse en la señal solar disponible en lugar del índice temporal genérico. Combinado con la ventaja sistemática, aunque de magnitud pequeña, de LayerNorm sobre Dynamic Tanh confirmada mediante un test de Diebold-Mariano con varianza Newey-West (11 de 12 combinaciones significativas a favor de LN), este hallazgo motiva la adopción de \textbf{S+C-LN} como modelo de referencia único para el resto del trabajo, prescindiendo de Prophet, del modelo de residuos y del ensamblaje por Stacking Optimization.

El segundo resultado central corresponde a Conformal Prediction. Sobre S+C-LN, se implementaron cuatro variantes que comparten el mismo predictor puntual crudo (Split CP clásico, ACI, LW Split CP y LW-ACP), más regresión cuantílica conformada (CQR) como comparación externa a la familia de CP, con $\gamma = 0.05$ fijado de antemano en los métodos adaptativos, en vez de seleccionado según su desempeño en el conjunto de prueba, calibrados excluyendo una ventana verificada de precio de falla administrativo (Sección \ref{sec:cp_configuracion_general}), y con el límite inferior de todos los intervalos truncado en 0 USD/MWh para respetar el dominio físico del costo marginal. El linaje conceptual Split CP $\to$ ACP $\to$ LW-ACP mostró que los dos mecanismos que lo componen son complementarios pero de contribución muy desigual: la adaptación \textit{online} del nivel de cobertura (ACI) por sí sola ya reduce la desviación de cobertura en un 57\% respecto al cuantil fijo \emph{y} el ancho del intervalo en un 8.0\%, mientras que la ponderación local por hora sola (LW Split CP) no reduce el ancho de forma consistente. La combinación de ambos mecanismos en LW-ACP añade una reducción adicional modesta pero no uniforme (angosta en cinco de las ocho barras), hasta un 8.7\% respecto al cuantil fijo --el más angosto de los cinco métodos, incluyendo CQR, que en esta comparación resulta ser el más ancho (67.99\,USD/MWh) y con una calibración notoriamente peor: 0.0081 de desviación de cobertura media, más del doble que ACI y LW-ACP--, con una cobertura empírica dentro de 0.24--0.60 puntos porcentuales del nominal en las ocho barras. LW-ACP es el único de los cuatro métodos de la familia conforme que combina el mejor resultado en ancho, calibración e \textit{interval score} de Winkler de forma simultánea. Esta separación limpia de efectos --posible únicamente al fijar un predictor puntual común entre métodos y al excluir la contaminación de calibración detectada-- muestra que la eficiencia de LW-ACP proviene mayoritariamente de la adaptación \textit{online}, con la heterocedasticidad horaria aportando un refinamiento adicional, no la fuente principal de la ganancia como se había estimado en un análisis preliminar de este mismo trabajo.

El análisis de errores extremos explica, a nivel mecanístico, por qué esta combinación es necesaria. En las ocho barras, entre el 80\% y el 100\% de los veinte mayores errores de validación corresponden a un \textit{desfase de fase}: el modelo predice correctamente la magnitud de la transición diaria del costo marginal, pero con un retraso sistemático de aproximadamente una hora respecto a la transición real, un patrón consistente en todo el SEN y no un artefacto de una barra particular. Estos errores son de gran amplitud pero muy localizados en el tiempo, lo que explica la brecha entre el MAE del modelo y el ancho de los intervalos de Split CP clásico, y motiva directamente el valor de las variantes adaptativas y ponderadas por hora.

Puerto Montt emergió de forma consistente, a lo largo de todo el trabajo, como la barra de peor desempeño del sistema, con un MAE de test de 13.68\,USD/MWh frente a 6.4--8.0\,USD/MWh del resto y un R² de 0.823, el más bajo del SEN. Las estadísticas descriptivas descartan que se trate de una barra ``difícil'' al azar: su media y desviación estándar de precio son sustancialmente mayores que las del resto, consistente con una dependencia hidrológica de baja frecuencia que opera en escalas de meses, más lentas que la ventana de contexto de 168 horas del TFT. Los experimentos con \textit{features} hídricas propias (H+C y S+H+C) no lograron mejorar sobre S+C para esta barra, y el análisis de errores extremos mostró que, aunque Puerto Montt también exhibe el patrón de desfase de fase, este no está ligado al ciclo solar como en el resto del SEN, sino a una dinámica de temporización propia. Puerto Montt queda documentado como un caso estructuralmente distinto, no resuelto por las extensiones exploradas en este trabajo.

El estudio de granularidad temporal, explorado sobre tres barras piloto (Atacama, Cardones y Puerto Montt) y evaluado contra la serie cruda sin recorte de outliers, sugiere que una mayor resolución temporal mejora la capacidad predictiva del modelo, aunque de forma más moderada de lo que sugería la serie recortada: el R² pasa de $\sim$0.27--0.45 a resolución diaria, a $\sim$0.82--0.86 a resolución horaria, hasta $\sim$0.85--0.88 a 15 minutos, con la mejora más marcada precisamente en Puerto Montt ($\sim$0.36--0.42 $\to$ $\sim$0.82 $\to$ $\sim$0.87). El salto de Día a Horario domina la mejora total; el paso adicional de Horario a 15 minutos aporta entre 2 y 5 puntos porcentuales de R² en cinco de las seis combinaciones de barra y arquitectura, y no mejora en la restante (Cardones LN), dentro del margen de ruido esperable --el recorte de outliers original exageraba esta última mejora al afectar proporcionalmente más a la resolución de 15 minutos--. La hipótesis se sostiene con más fuerza al controlar por el \textit{confound} de período histórico: comparando Horario y 15 minutos sobre exactamente la misma ventana calendario, un resultado que no se vio afectado por la corrección de outliers, el MAE puntual se reduce entre un 44\% y un 55\% al aumentar la resolución (por ejemplo, Puerto Montt: 3.58 vs.\ 7.87\,USD/MWh), lo que descarta que la ventaja observada sea enteramente un artefacto del período más reciente cubierto por los datos de mayor resolución. La comparación con la granularidad Día, en cambio, no pudo controlarse de la misma forma y queda como \textit{confound} abierto; extender este control y el estudio a las ocho barras es la continuación natural de este resultado.

Finalmente, la comparación con literatura reciente (León et al., 2026), que evalúa modelos de \textit{machine learning} clásico sobre barras compartidas del SEN, mostró que el TFT S+C obtiene un MAE 34\% menor en Crucero (5.11 vs.\ 7.71\,USD/MWh) y un 20\% menor en Puerto Montt (13.64 vs.\ 17.11\,USD/MWh) respecto a los mejores modelos reportados, pese a entrenar con menos de un año de datos en ambos casos y sin el ajuste específico por barra que caracteriza al pipeline principal de este trabajo. En Puerto Montt, el R² resultó notoriamente más bajo que el de León et al. (0.56 vs.\ 0.76), una aparente contradicción que se explica por la menor varianza del precio en la ventana de evaluación de este trabajo: relativizando el error cuadrático medio contra la varianza de la ventana de León et al.\ se obtiene un R² implícito de $\approx 0.83$, consistente con el $0.823$ que Puerto Montt alcanza en el pipeline principal, lo que confirma que el MAE y el RMSE, métricas que no dependen de la varianza del conjunto evaluado, son las más directamente comparables entre ambos trabajos.

\section{Contribuciones}

Las contribuciones principales de este trabajo son:

\begin{itemize}
    \item \textbf{LW-ACP (\textit{Locally-Weighted Adaptive Conformal Prediction})}: la contribución metodológica central del trabajo, que combina la ponderación local por hora de LW Split CP con la adaptación \textit{online} del nivel de cobertura de ACI (Gibbs y Candès, 2021), demostrando empíricamente que ambos mecanismos son complementarios --aunque de magnitud desigual, con la adaptación \textit{online} explicando la mayor parte de la ganancia-- y no redundantes para series de costo marginal eléctrico con heterocedasticidad horaria pronunciada.
    \item Evidencia sistemática, respaldada por un test de significancia estadística y no solo por métricas agregadas, de que reemplazar la descomposición de Prophet por la posición solar como \textit{feature} exógena simplifica el modelo sin sacrificar precisión, y en la mayoría de los casos la mejora.
    \item Una caracterización cuantitativa y reproducible del mecanismo detrás de los errores extremos del modelo, mediante un criterio de desplazamiento temporal que distingue errores de \textit{desfase de fase} de \textit{eventos estructurales}, aplicable en principio a otros modelos de pronóstico de series horarias con ciclos diarios marcados.
    \item Un tratamiento diferenciado y documentado de Puerto Montt, que identifica sus features hídricas como una hipótesis razonable pero no confirmada, evitando la conclusión apresurada de que un ajuste de \textit{features} resuelve su desempeño estructuralmente inferior.
    \item Evidencia, verificada controlando el \textit{confound} de período histórico entre granularidades, de que una mayor resolución temporal (15 minutos vs.\ horario) mejora sistemáticamente la capacidad predictiva del modelo, una hipótesis con implicancias directas para la operación de mercados eléctricos que despachan a intervalos sub-horarios y que motiva extender el pipeline principal de este trabajo más allá de la resolución horaria.
    \item Una comparación aplicada, ausente del diseño original de este trabajo, contra el costo marginal programado que el CEN publica operacionalmente: en las seis barras donde se logró identificar el mnemotécnico correspondiente, S+C-LN reduce el MAE entre 52.1\% y 60.6\% respecto a esa cifra institucional (promedio 55.5\%, incluyendo Puerto Montt, la barra de peor desempeño del sistema en el resto de este trabajo; Sección \ref{sec:cen_programado}), evidencia de que un modelo dedicado y sin las restricciones operativas de la Programación de la Operación puede superar de forma sistemática, no solo puntual, al valor que efectivamente circula en el mercado antes del despacho real.
    \item Una evaluación de sensibilidad frente a una función de pérdida asimétrica (Sección \ref{sec:perdida_asimetrica}), ausente de las métricas simétricas usadas en el resto de este trabajo: la ventaja de S+C-LN sobre la persistencia se sostiene en siete de las ocho barras independientemente de qué tipo de error --sobreestimar o subestimar el precio-- resulte más costoso para el agente, pero en Puerto Montt esa conclusión sí depende de la asimetría asumida, un matiz que el MAE no puede revelar por construcción y que resulta relevante para cualquier lectura aplicada de estos resultados.
\end{itemize}

\section{Limitaciones}

El esquema de calibración de Conformal Prediction presenta una limitación metodológica inherente al diseño experimental: el conjunto de validación cumple simultáneamente el rol de selección del punto de \textit{early-stopping} del TFT y de calibración del cuantil de no conformidad, lo que introduce un sesgo optimista esperablemente pequeño dado el tamaño del conjunto de calibración ($n_\text{cal} = 8\,298$), pero presente.

EnbPI no fue implementado sobre S+C: su mecanismo de \textit{bootstrap} está diseñado para un modelo base entrenable con múltiples réplicas, mientras que S+C es una única predicción del TFT sin ensamble subyacente que bootstrapear. ACI, en cambio, sí se adaptó a S+C directamente, al operar sobre el mismo \textit{score} de no conformidad que Split CP en lugar de requerir un ensamble; EnbPI habría necesitado una reformulación más profunda de su procedimiento para prescindir del \textit{bootstrap}.

El estudio de granularidad temporal es exploratorio: se restringió a tres barras piloto. En el proceso reveló una inestabilidad del estimador $\hat\sigma_h$ de LW Split CP y LW-ACP cuando el conjunto de calibración a 15 minutos queda estacionalmente sesgado y algunas horas presentan varianza casi nula --el mismo mecanismo, agravado por una ventana de calibración corta, que motivó el piso de varianza $\hat\sigma_h \geq \hat\sigma_{\text{global}}$ adoptado en el pipeline principal (Sección \ref{sec:lw_split_cp})--, quedando resuelta con esa misma corrección. La comparación entre granularidades Día y Horario tampoco pudo controlarse completamente por el confound de período histórico, a diferencia de la comparación Horario vs. 15 minutos.

Todo el pipeline, no solo el Stacking Optimization (Sección 2.5.2), se evalúa sobre un único split cronológico train/validación/prueba, en vez de \textit{walk-forward validation} sobre múltiples ventanas rodantes. Esto implica que las cifras del Capítulo 4 provienen de un solo período de prueba (mayo 2025--abril 2026), sin una estimación de cuánto varía el desempeño del modelo si ese período se desplaza --una práctica recomendada específicamente para forecasting de precios eléctricos dado que el SEN es un mercado en transición estructural continua (Lago et al., 2021)--. La razón es de costo computacional: entrenar un único $\text{TFT}_\text{LN}$ toma 3--3.5 horas en GPU, y una época completa en CPU toma 5.96 minutos, lo que extrapolado a las 100 épocas que entrena un modelo como el de Atacama implicaría cerca de 10 horas por modelo en CPU; repetir esto sobre varias ventanas rodantes, ocho barras y dos arquitecturas sin acceso a cluster no era viable dentro del plazo de este trabajo.

El esfuerzo de ajuste de hiperparámetros tampoco es simétrico entre modelos: Prophet se optimiza mediante una grilla de \texttt{changepoint\_prior\_scale} y \texttt{seasonality\_prior\_scale} específica para cada barra, mientras que los hiperparámetros del TFT se fijan de antemano para las ocho barras y ambas arquitecturas, dado el costo de entrenar cada TFT (horas, frente a segundos por combinación de Prophet). Esto introduce una limitación explícita sobre dos comparaciones de este trabajo: la ventaja de S+C sobre Prophet solo podría, en principio, ampliarse o reducirse con hiperparámetros del TFT optimizados por barra; y la ventaja de LayerNorm sobre DyT se mide bajo una configuración no optimizada específicamente para ninguna de las dos arquitecturas, por lo que no puede descartarse que una búsqueda independiente para cada una altere esa conclusión.

El pipeline tampoco es directamente desplegable en operación real tal como está evaluado: las variables climáticas provienen de ERA5, un producto de reanálisis que se publica con días de rezago, por lo que ni siquiera las horas más recientes de la ventana de encoder estarían disponibles al momento de generar una predicción en tiempo real; lo mismo aplica, en menor escala, a las variables hídricas de Puerto Montt (H+C, S+H+C), cuyo valor diario se expande a las 24 horas del día sin que ese valor esté necesariamente cerrado desde las 00:00. El esquema es causalmente correcto para el entrenamiento y la evaluación de este trabajo, pero una implementación operacional requeriría sustituir ambas fuentes por productos de pronóstico o por el último valor cerrado disponible con rezago.

Finalmente, Puerto Montt permanece como un caso no resuelto: ninguna de las estrategias evaluadas en este trabajo, incluyendo las \textit{features} hídricas diseñadas específicamente para esta barra, logra llevar su desempeño a la par del resto del SEN.

\section{Trabajo futuro}

A partir de las limitaciones anteriores y de los hallazgos del trabajo, se identifican las siguientes líneas de continuación:

\begin{itemize}
    \item Extender EnbPI a S+C mediante un esquema de \textit{bootstrap} directo sobre el TFT, para completar la comparación entre las tres familias de Conformal Prediction adaptativo bajo un mismo pipeline y el mismo predictor puntual.
    \item Extender la comparación contra el costo marginal programado del CEN (Sección \ref{sec:cen_programado}) a Atacama y Tarapacá, las dos barras restantes del SEN, cuyo mnemotécnico no se logró ubicar en el catálogo de la API pública del CEN pese a intentarlo con las variantes de voltaje y sección de barra disponibles --la consulta a la API y el cálculo de la comparación ya están automatizados para el resto--.
    \item Evaluar el pipeline mediante \textit{walk-forward validation} sobre al menos dos o tres ventanas rodantes adicionales, para cuantificar cuánto varían las cifras del Capítulo 4 según el período de prueba, en vez de reportar un único split. El costo estimado (Sección Limitaciones) es de aproximadamente 3--3.5 horas por modelo en GPU, lo que hace viable repetirlo para un subconjunto de barras si se dispone de acceso a cluster.
    \item Investigar directamente el mecanismo del desfase de fase identificado en el análisis de errores extremos, por ejemplo mediante una función de pérdida que penalice explícitamente el error de temporización, o una corrección \textit{post-hoc} del desplazamiento sistemático de aproximadamente una hora observado en las ocho barras.
    \item Extender el estudio de granularidad temporal a las ocho barras del SEN y a un período de calibración de 15 minutos que cubra un año completo, para resolver el confound estacional identificado --la inestabilidad de $\hat\sigma_h$ en sí ya quedó resuelta con el piso de varianza adoptado (Sección \ref{sec:lw_split_cp}), independiente de si la ventana de calibración es estacionalmente balanceada--.
    \item Explorar un ensamble entre $\text{TFT}_\text{LN}$ y $\text{TFT}_\text{DyT}$ mediante Stacking Optimization, dado que ambas arquitecturas exhiben patrones de error suficientemente distintos como para aportar diversidad complementaria, pese a que LN domina de forma sistemática cuando se evalúa individualmente.
    \item Evaluar la aplicabilidad de LW-ACP a otros sistemas eléctricos con alta penetración de generación solar, donde la heterocedasticidad horaria asociada al ciclo diario de irradiancia sería, a priori, un patrón igualmente explotable.
    \item Explorar una variante de LW-ACP con adaptación \textit{online} indexada por hora del día (24 trayectorias de $\alpha_t$ independientes en lugar de una única trayectoria global), de modo que la velocidad de reacción ante aciertos y fallos ($\gamma$) se ajuste específicamente en torno a las horas de transición solar sin arrastrar el ensanchamiento a las horas adyacentes, mitigando el mismo efecto de desfase temporal identificado en las predicciones puntuales. Esta variante ya no es solo una hipótesis: la cobertura condicional por hora (Sección \ref{sec:cobertura_condicional}) confirma que el $\alpha_t$ único y compartido de LW-ACP sub-cubre específicamente en las horas de transición solar (0.898 frente al 0.95 nominal), mientras que su cobertura mes a mes es notablemente estable, evidencia directa de que el problema es la falta de granularidad horaria del mecanismo \textit{online}, no de la adaptación \textit{online} en sí misma.
\end{itemize}
